%  LaTeX support: latex@mdpi.com 
%  For support, please attach all files needed for compiling as well as the log file, and specify your operating system, LaTeX version, and LaTeX editor.

%=================================================================
\documentclass[atmosphere,article,submit,pdftex,moreauthors]{Definitions/mdpi} 

%=================================================================
% MDPI internal commands
\firstpage{1} 
\makeatletter 
\setcounter{page}{\@firstpage} 
\makeatother
\pubvolume{1}
\issuenum{1}
\articlenumber{0}
\pubyear{2022}
\copyrightyear{2022}
%\externaleditor{Academic Editor: Firstname Lastname}
\datereceived{} 
%\daterevised{} % Only for the journal Acoustics
\dateaccepted{} 
\datepublished{} 
%\datecorrected{} % Corrected papers include a "Corrected: XXX" date in the original paper.
%\dateretracted{} % Corrected papers include a "Retracted: XXX" date in the original paper.
\hreflink{https://doi.org/} % If needed use \linebreak
%\doinum{}
%------------------------------------------------------------------
% The following line should be uncommented if the LaTeX file is uploaded to arXiv.org
%\pdfoutput=1

%=================================================================
% Add packages and commands here. The following packages are loaded in our class file: fontenc, inputenc, calc, indentfirst, fancyhdr, graphicx, epstopdf, lastpage, ifthen, lineno, float, amsmath, setspace, enumitem, mathpazo, booktabs, titlesec, etoolbox, tabto, xcolor, soul, multirow, microtype, tikz, totcount, changepage, attrib, upgreek, cleveref, amsthm, hyphenat, natbib, hyperref, footmisc, url, geometry, newfloat, caption
\graphicspath{{figures/}} % path to folder with figures
\usepackage{subcaption}
\usepackage{listings}
\usepackage{xcolor}
\usepackage{gensymb} % for \textdegree
\usepackage{tabularx}

%=================================================================
%% Please use the following mathematics environments: Theorem, Lemma, Corollary, Proposition, Characterization, Property, Problem, Example, ExamplesandDefinitions, Hypothesis, Remark, Definition, Notation, Assumption
%% For proofs, please use the proof environment (the amsthm package is loaded by the MDPI class).

%=================================================================
% Full title of the paper (Capitalized)
\Title{Mapping Land Cover Changes and Dynamics of Deforestation in Burkina Faso Based on Landsat Scenes and TerraClimate Data}

% MDPI internal command: Title for citation in the left column
\TitleCitation{Mapping Land Cover Changes and Dynamics of Deforestation in Burkina Faso Based on Landsat Scenes and TerraClimate Data}

% Author Orchid ID: enter ID or remove command
\newcommand{\orcidauthorA}{0000-0002-5759-1089} % Add \orcidA{} behind the author's name
\newcommand{\orcidauthorB}{0000-0002-6461-1551} % Add \orcidB{} behind the author's name

% Authors, for the paper (add full first names)
\Author{Polina Lemenkova $^{1}$\orcidA{} and Olivier Debeir $^{1}\orcidB{}$}
% \dagger,\ddagger
%\longauthorlist{yes}

% MDPI internal command: Authors, for metadata in PDF
\AuthorNames{Polina Lemenkova $^{1}$ and Olivier Debeir $^{1}$}

% MDPI internal command: Authors, for citation in the left column
\AuthorCitation{Lemenkova, P.; Debeir, O.}
% If this is a Chicago style journal: Lastname, Firstname, Firstname Lastname, and Firstname Lastname.

% Affiliations / Addresses (Add [1] after \address if there is only one affiliation.)
\address{%
$^{1}$ \quad Laboratory of Image Synthesis and Analysis (LISA), École polytechnique de Bruxelles (Brussels Faculty of Engineering), Université Libre de Bruxelles (ULB). Building L, Campus du Solbosch, ULB -- LISA CP165/57, Avenue Franklin D. Roosevelt 50, B-1050, Brussels, Belgium}

% Contact information of the corresponding author
\corres{Correspondence: polina.lemenkova@ulb.be; Tel.: +32471860459 (P.L.)}

% Abstract (Do not insert blank lines, i.e. \\) 
\abstract{
Situated in the southern Sahel zone, Burkina Faso experiences one of the most extreme climatic hazards in Sub-Saharan Africa. These vary from extreme floods in Volta River Basin, one of the most vulnerable regions in West Africa to hydrological disasters, to desertification and recurrent droughts caused by the rise in annual temperatures and decreased runoff. Besides climate factors, anthropogenic activities contribute to land degradation and deforestation through agricultural lands expansion of croplands, which leads to the decrease of woody savannah and gallery forests in the country. High exposure and low adaptive capacity to climate-related hazards require development of adaptive strategies in Burkina Faso with aim to control food management, secure agricultural production and maintain environmental sustainability. In this paper, we demonstrated the integrated use of the remote sensing data and climate parameters for mapping climate parameters and vegetation dynamics using scripting cartographic techniques. For problems of computer vision and analytical modelling in cartography, reliable datasets are essential. Therefore, we used the satellite imagery scenes from the Landsat 8-9 OLI/TIRS C2 L1 collection for mapping changes in forest and savannah coverage in Burkina Faso, and TerraClimate data for visualizing key climate parameters. Using integrated dataset we visualised the dynamics of target climate characteristics of Burkina Faso for 2013-2022: extreme temperatures, level of precipitation, soil moisture, downward surface shortwave radiation, vapor pressure deficit and anomaly. Additionally, the Palmer Drought Severity Index (PDSI) was modelled over the study area to estimate soil water balance related to soil moisture conditions as a prerequisites for vegetation growth. This paper contributes the to environmental analysis and monitoring of Burkina Faso with aim to highlight the links between climate processes and vegetation dynamics in West Africa. 
}

% Keywords
\keyword{keyword 1; keyword 2; keyword 3 (List three to ten pertinent keywords specific to the article; yet reasonably common within the subject discipline.)} 

% The fields PACS, MSC, and JEL may be left empty or commented out if not applicable
\PACS{91.10.Da; 91.10.Jf; 91.10.Sp; 91.10.Xa; 96.25.Vt; 91.10.Fc; 95.40.+s; 95.75.Qr; 95.75.Rs; 42.68.Wt}
\MSC{86A30; 86-08; 86A99; 86A04}
\JEL{Y91; Q20; Q24; Q23; Q3; Q01; R11; O44; O13; Q5; Q51; Q55; N57; C6; C61}

%%%%%%%%%%%%%%%%%%%%%%%%%%%%%%%%%%%%%%%%%%
\begin{document}

%%%%%%%%%%%%%%%%%%%%%%%%%%%%%%%%%%%%%%%%%%
\section{Introduction}

\subsection{Background}
Despite years of active research, the role of climate in vegetation remains an unsolved problem for a number of reasons, First, the response of vegetation to climate has a nonlinear character and complex nature with both direct and indirect effects. For instance, a variety of physical geographic and climatic factors directly contribute to the development of the vegetation. These include temperature and solar radiation, moisture and precipitation, wind and relief morphometry. The indirect processes include the effects from the climatic parameters on soil conditions, health of microorganisms and fauna that constitute a vital part of the complex biochemical cycle and food chain. 

Second, revealing complex interactions between vegetation and climate requires the analyses of large datasets to highlight links and correlations between the parameters. This necessarily requires advanced methods of big data processing. Moreover, even moderate changes in climate variables  often result in a noticeable change in vegetation patters in a spatio-temporal perspective. These can be analysed using technical methods of predictive modelling and simulation of the spatial distribution of the vegetation patterns across a landscape based on climatic and environmental variables as input data. All of these factors significantly alter our understanding of how climate-vegetation links function.

Third, the machine-based quantification of complex climate-vegetation system requires considering not only a full range of climatic and environmental parameters that affect vegetation condition and growth but also a temporal delay in reaction of plant species on climate change. A complete change in land cover types and landscape patterns, including boundaries of habitats, may take up to several seasons. Further, simulating climate-environmental relationships also needs to handle other external factors such as urban growth and industrial heat sources, water pollution and soil contaminations, as well as changes in local topography which may have a human-induced (construction works, career and mining) and natural origin (landslides, soil subsidence).

While all of these factors have been studied, the problem of advanced methods of data processing is the only one that we believe is still largely unsolved. Most commercially deployed Geographic Information System (GIS) are based on the predefined functionality of the software and have limited embedded algorithms of data processing. They sometimes employ simple methods of image manipulation which does not take into consideration real features of datasets. Processing large geospatial datasets becomes even more critical for standard GIS tools that rely on manual data processing. In contrast, the advanced scripting solutions and  programming tools enable to handle geospatial data automatically and accurately, which is the main focus on this paper.

This paper is structured as follows. Section 1 overviews the existing recent efforts made to tackle climate-environmental relationships in recent literature. Section 2 gives a brief overview of the study area focused on Burkina Faso, French West Africa. Here we discussed the major geographic features this region, environmental challenges and the effects from recent climate changes on agricultural system of the country. In Section 3 we argue that we need an essential advanced mapping tools for accurate and automated data visualization. For that we apply a fast and efficient scripting methods of the Generic Mapping Toolset (GMT) for mapping climate-environmental parameters. Libraries of R language are applied using programming techniques to improve the workflow of satellite image processing. Section 4 details the obtained results and observed variations in environmental parameters in Burkina Faso at regional scale and at local scale with enlarged fragment of the study area. Section 5 analyses the robustness of scripting techniques with respect to mapping climate and environmental parameters and performed image processing, including computed NDVI and visualized land cover changes in Burkina Faso over the years 2013-2022.

\subsection{Related Work}


\subsection{Contribution}
In this work we consider the problem of satellite image processing for climate and environmental mapping using programming methods and machine learning approaches. Over the last few years this topic has received a growing amount of attention in the Earth sciences and geospatial community \cite{4779429,jimaging8120317,DIMOBE2022100768,app122412554}. We argue, however, that nearly all existing research papers on environmental mapping of Botswana have focused on using traditional Geographic Information Systems (GIS) \cite{Augusseau,AFR350,hydrology9010012,land7040124,ZOUNGRANA201866,Forkuor}. Such approaches reflected the convenient functionality and characteristics of the most popular GIS software with Graphical User's Interface (GUI) such as Erdas Imaging, ArcGIS, ENVI GIS, Idrisi GIS, SAGA GIS and others. Nevertheless, the programming approaches offer more advanced performance of spatial data processing in terms of computational effectiveness, speed and effectiveness of mapping, image classification accuracy and oprimised cartographic workflow, due to the high level of automation and repeatability of scripts.

\section{Study Area}

The problems of climate change in sub-Saharan Africa are largely discussed in many recent papers \cite{AYANLADE2022100311,EMEDIEGWU2022105967,WAONGO201523,AZIBO2015126,ASCOTT2022128107,WEBBER2014161,WAHA2013130}. Major consequences of climate change include the decline in precipitation and increase in annual temperatures. Such trends results in changed vegetation patterns, lead to deforestation and desertification in West Africa, and negatively affect food production in agricultural sector, including both rainfed and irrigated crop yields \cite{CALZADILLA2013150,WAONGO201523}. 

One of the sub-Saharan Africa countries situated in the southern Sahel zone (Figure \ref{fig01}), Burkina Faso experiences extreme climatic hazards \cite{Zougmore}. These vary from the extreme floods in Volta River Basin, one of the most vulnerable regions in West Africa to hydrological disasters, desertification and recurrent droughts caused by the rise in annual temperatures and decreased runoff. As a result, this poses high risks to the vegetation sustainability, food security in mixed crop–livestock systems \cite{RIGOLOT2017217}. Other reported consequences of climate impacts in Burkina Faso include changes in land cover types with detected increase in tree savannah, bare soils and agricultural lands, and decrease in woodland, gallery forest, shrub savannah and water areas \cite{DIMOBE2015559}. Besides climate factors, anthropogenic activities contribute to land degradation and deforestation through agricultural lands expansion of croplands, which leads to the decrease of woody savannah and gallery forests in the country. 

\begin{figure}[H]
\centering
	\includegraphics[width=14.5 cm]{Fig_01.jpg}
	\caption{Topographic map of Burkina Faso, West Africa. Mapping software: GMT scripting toolset. Data source: GEBCO/SRTM. Cartography: authors.
	\label{fig01}}
\end{figure}

At the same time, the economy of Burkina Faso largely relies on agriculture and the use of natural resources as one of the least developed countries in the world with highest level of poverty \cite{1525946}. The common agricultural species are mostly planted in semi-urban areas and include, for instance, tomates, lettuce, cabbage, beans, eggplants \cite{Ouedraogo}. A moderate production of cotton (less than 1 t/ha) contributed to the development of traditional agrarian cropping systems in the savannahs of Burkina Faso and remains a notable feature of its agricultural and development policies \cite{Soumare}. However, recent issues of climate change resulted in continuous degradation in the agricultural-pastoral potentialities of the country, caused by both the intensive use of resources \cite{Vall} and climate change \cite{olsson1993desertification}. 

High exposure and low adaptive capacity to climate-related hazards require development of adaptive strategies in Burkina Faso with aim to control food management, secure agricultural production and maintain environmental sustainability \cite{BUNCLARK2018243,daniel2013farming}. Since the country is located in tropical climate, is has two distinct seasons - dry and wet - that control vegetation structure and the periods of agricultural works \cite{brons2000climate}. The increased drought periods and intensity, as well as intra-seasonal rainfall variability threaten to crop systems and livelihoods of Burkinabè farmers and pose high risks to sustainable yield harvesting \cite{su142315637,salter1967crop,atmos13071155}. In turn, the decrease of precipitation directly affect soil structure and biochemistry of plants. For instance, soil salinity increasing in periods of drought negatively affects the yields of maize, rice and other grain crops in Burkina Faso. Other consequences of climate change report variations in catchment hydrology and increased soil erosion, changes in evapotranspiration, water discharge and pattern of sedimentation \cite{OPDEHIPT2019431,OPDEHIPT201863}.

The vegetation of the Burkina Faso has a Sahelian type and is distributed according to the major topographic units with a latitudinal orientation of the Pleistocene dune belts \cite{Ballouche}. Geographically, the country consists of two major types of landscapes with varied vegetation types according to the topography of the country which is presented by a flat relief and gentle variations between the elevations \cite{Nebie}. The major part of the country has a gently undulating landscape and occasional isolated hills, while the contrasting southwest includes the highest peak, Ténakourou, on a sandstone plateau massif bordered by cliffs close to the border with Mali, Figure \ref{fig01}.

%%%%%%%%%%%%%%%%%%%%%%%%%%%%%%%%%%%%%%%%%%
\pagebreak
\section{Materials and Methods}

%%%%%%%%%%%%%%%%%%%%%%%%%%%%%%%%%%%%%%%%%%
\pagebreak
\section{Results}

\begin{figure}[H]
\makebox[0.90\linewidth][r]{%
     \begin{subfigure}[r]{0.5\textwidth}
         \caption{} \vspace{1mm}
         	\includegraphics[width=8.5cm]{BF_Tmin_2013.jpg}
     \end{subfigure}\hspace{12mm}
     \begin{subfigure}[r]{0.5\textwidth}
             \caption{} \vspace{1mm}
         	\includegraphics[width=8.5cm]{BF_Tmin_2020.jpg}
     \end{subfigure}
   }
  \caption{Minimum temperature for Burkina Faso. \textbf{(a)}: 2013; \textbf{(b)}: 2020.}
        \label{figTmin}
\end{figure}

\begin{figure}[H]
\makebox[0.90\linewidth][r]{%
     \begin{subfigure}[r]{0.5\textwidth}
         \caption{} \vspace{1mm}
         	\includegraphics[width=8.5cm]{BF_Tmax_2013.jpg}
     \end{subfigure}\hspace{12mm}
     \begin{subfigure}[r]{0.5\textwidth}
             \caption{} \vspace{1mm}
         	\includegraphics[width=8.5cm]{BF_Tmax_2020.jpg}
     \end{subfigure}
   }
  \caption{Maximum temperature for Burkina Faso. \textbf{(a)}: 2013; \textbf{(b)}: 2020.}
        \label{figTmax}
\end{figure}

\begin{figure}[H]
\makebox[0.90\linewidth][r]{%
     \begin{subfigure}[r]{0.5\textwidth}
         \caption{} \vspace{1mm}
         	\includegraphics[width=8.5cm]{BF_soil_2013.jpg}
     \end{subfigure}\hspace{12mm}
     \begin{subfigure}[r]{0.5\textwidth}
             \caption{} \vspace{1mm}
         	\includegraphics[width=8.5cm]{BF_soil_2020.jpg}
     \end{subfigure}
   }
  \caption{Soil moisture in Burkina Faso. \textbf{(a)}: 2013; \textbf{(b)}: 2020.}
        \label{figSoil}
\end{figure}

\begin{figure}[H]
\makebox[0.90\linewidth][r]{%
     \begin{subfigure}[r]{0.50\textwidth}
         \caption{} \vspace{1mm}
         	\includegraphics[width=8.5cm]{BF_pet_2013.jpg}
     \end{subfigure}\hspace{12mm}
     \begin{subfigure}[r]{0.50\textwidth}
             \caption{} \vspace{1mm}
         	\includegraphics[width=8.5cm]{BF_pet_2020.jpg}
     \end{subfigure}
}
  \caption{Potential evapotranspiration (PET) in Burkina Faso. \textbf{(a)}: 2013; \textbf{(b)}: 2020.}
        \label{figPET}
\end{figure}

\begin{figure}[H]
\makebox[0.90\linewidth][r]{%
     \begin{subfigure}[r]{0.50\textwidth}
         \caption{} \vspace{1mm}
         	\includegraphics[width=8.5cm]{BF_aet_2013.jpg}
     \end{subfigure}\hspace{12mm}
     \begin{subfigure}[r]{0.50\textwidth}
             \caption{} \vspace{1mm}
         	\includegraphics[width=8.5cm]{BF_aet_2020.jpg}
     \end{subfigure}
}
  \caption{Actual evapotranspiration (AET) in Burkina Faso. \textbf{(a)}: 2013; \textbf{(b)}: 2020.}
        \label{figAET}
\end{figure}

\begin{figure}[H]
\makebox[0.90\linewidth][r]{%
     \begin{subfigure}[r]{0.50\textwidth}
         \caption{} \vspace{1mm}
         	\includegraphics[width=8.5cm]{BF_VPD_2013.jpg}
     \end{subfigure}\hspace{12mm}
     \begin{subfigure}[r]{0.50\textwidth}
             \caption{} \vspace{1mm}
         	\includegraphics[width=8.5cm]{BF_VPD_2020.jpg}
     \end{subfigure}
}
  \caption{Vapour Pressure Deficit (VPD) in Burkina Faso. \textbf{(a)}: 2013; \textbf{(b)}: 2020.}
        \label{figVPD}
\end{figure}

\begin{figure}[H]
\makebox[0.90\linewidth][r]{%
     \begin{subfigure}[r]{0.50\textwidth}
         \caption{} \vspace{1mm}
         	\includegraphics[width=8.5cm]{BF_VAP_2013.jpg}
     \end{subfigure}\hspace{12mm}
     \begin{subfigure}[r]{0.50\textwidth}
             \caption{} \vspace{1mm}
         	\includegraphics[width=8.5cm]{BF_VAP_2020.jpg}
     \end{subfigure}
}
  \caption{Vapour Pressure (VAP) in Burkina Faso. \textbf{(a)}: 2013; \textbf{(b)}: 2020.}
        \label{figVAP}
\end{figure}

\begin{figure}[H]
\makebox[0.90\linewidth][r]{%
     \begin{subfigure}[r]{0.50\textwidth}
         \caption{} \vspace{1mm}
         	\includegraphics[width=8.5cm]{BF_SRAD_2013.jpg}
     \end{subfigure}\hspace{12mm}
     \begin{subfigure}[r]{0.50\textwidth}
             \caption{} \vspace{1mm}
         	\includegraphics[width=8.5cm]{BF_SRAD_2020.jpg}
     \end{subfigure}
}
  \caption{Downward Surface Shortwave Radiation (SRAD) in Burkina Faso. \textbf{(a)}: 2013; \textbf{(b)}: 2020.}
        \label{figAET}
\end{figure}

\begin{figure}[H]
\makebox[0.90\linewidth][r]{%
     \begin{subfigure}[r]{0.50\textwidth}
         \caption{} \vspace{1mm}
         	\includegraphics[width=8.5cm]{BF_ws_2013.jpg}
     \end{subfigure}\hspace{12mm}
     \begin{subfigure}[r]{0.50\textwidth}
             \caption{} \vspace{1mm}
         	\includegraphics[width=8.5cm]{BF_ws_2020.jpg}
     \end{subfigure}
}
  \caption{Average wind speed in Burkina Faso. \textbf{(a)}: 2013; \textbf{(b)}: 2020.}
        \label{figWS}
\end{figure}

\begin{figure}[H]
\makebox[0.90\linewidth][r]{%
     \begin{subfigure}[r]{0.50\textwidth}
         \caption{} \vspace{1mm}
         	\includegraphics[width=8.5cm]{BF_PDSI_2013.jpg}
     \end{subfigure}\hspace{12mm}
     \begin{subfigure}[r]{0.50\textwidth}
             \caption{} \vspace{1mm}
         	\includegraphics[width=8.5cm]{BF_PDSI_2020.jpg}
     \end{subfigure}
}
  \caption{PDSI (Palmer Drought Severity Index) in Burkina Faso. \textbf{(a)}: 2013; \textbf{(b)}: 2020.}
        \label{figPDSI}
\end{figure}

%%%%%%%%%%%%%%%%%%%%%%%%%%%%%%%%%%%%%%%%%%
\section{Discussion}

%%%%%%%%%%%%%%%%%%%%%%%%%%%%%%%%%%%%%%%%%%
\section{Conclusions}

To handle the complex environmental scenarios encountered in the Burkina Faso due to the climate change, we propose a pairwise comparative mapping and image processing of the data on years 2013 and 2020 for each evaluated climatic and environmental parameter. The presented data visualization by scripts is a semi-automated mapping process run trough a console which is interleaved with adjustment of selected parameters in script for each map according to the data range and extend. Using scripts, we are able to incorporate the spatial contiguity of the geospatial continuum fields and analyse their variability over space, thereby fully realising the potential of the scripting cartographic method. The geospatial representation approach offered by the GMT scripting toolset and R language returns a series of maps showing various environmental parameters over the area of Burkina Faso. Extensive experiments with GMT and R scripts on challenging TerraClimate dataset demonstrate that our method outperforms various state-of-the-art GIS approaches and can be extended on a wider range of complex environmental scenarios in West Africa.

%%%%%%%%%%%%%%%%%%%%%%%%%%%%%%%%%%%%%%%%%%
\vspace{6pt} 

\authorcontributions{Supervision, conceptualization, methodology, software, resources, funding acquisition, and project administration, O.D.; writing---original draft preparation, methodology, software, data curation, visualization, formal analysis, validation, writing---review and editing, and investigation, P.L. All authors have read and agreed to the published version of the manuscript.}

\funding{The publication was funded by the Editorial Office of XXXX, Multidisciplinary Digital Publishing Institute (MDPI), by providing 100\% discount for the APC of this manuscript. This project was supported by the Federal Public Planning Service Science Policy or Belgian Science Policy Office, Federal Science Policy - BELSPO (B2/202/P2/SEISMOSTORM).}

\institutionalreview{Not applicable}

\informedconsent{Not applicable}

\dataavailability{Not applicable} 

\acknowledgments{The authors thank the anonymous reviewers of Land for reading, suggestions and comments that improved an earlier version of this manuscript.}

\conflictsofinterest{The authors declare no conflict of interest} 

\abbreviations{Abbreviations}{
The following abbreviations are used in this manuscript:\\

\noindent 
\begin{tabular}{@{}ll}
MDPI & Multidisciplinary Digital Publishing Institute\\
DOAJ & Directory of open access journals\\
TLA & Three letter acronym\\
LD & Linear dichroism
\end{tabular}
}

\reftitle{References}

%=====================================
% References, variant A: external bibliography
%=====================================
\bibliography{BurkinaFaso.bib}
%\nocite{*}

\end{document}

